\documentclass[11pt,a4paper]{article}
\usepackage[margin=1in]{geometry}
\usepackage{hyperref}
\usepackage{enumitem}
\usepackage{graphicx}
\usepackage{xcolor}

% -----------------------------------------------------
% bvraghav-tiet-ucs503p-article.sty
% -----------------------------------------------------

% Variable: Lab Instructor
\makeatletter \newcommand{\BvrSetLabInstructor}[1]{%
  \gdef\Bvr@LabInstructor{#1}}
% default
\newcommand{\Bvr@LabInstructor}{Dr. Raghav %
  B. Venkataramaiyer}
% public accessor
\newcommand{\BvrLabInstructorValue}{\Bvr@LabInstructor}
\makeatother

% Add subtitle
\usepackage{titling}
% -- Set title bold
\pretitle{\begin{center}\bfseries\fontsize{18bp}{18bp}\selectfont}
\makeatletter
\newcommand{\BvrSubtitle}[1]{%
  \posttitle{%
    \par\end{center}
    \begin{center}\large#1\end{center}
    \vskip 0.5em}%
}
\makeatother

% Hints in the section title(s)
\newcommand{\BvrHint}[1]{%
  {\normalfont\normalsize\itshape\hfill #1}}

% -----------------------------------------------------

\title{Vyapari: A Unified ERP Platform for Retail and Wholesale Businesses}

\BvrSetLabInstructor{Miss Paramveer Kaur}

\BvrSubtitle{%
  Project Proposal for UCS503P  \\
  Submitted to: \BvrLabInstructorValue \\ \bfseries
  Thapar Institute of Engineering and Technology% 
}

\author{
  Vanshika, CSED%
  \thanks{Roll No: \texttt{1024030163}, %
    Email: \texttt{vvanshika\_be24@thapar.edu}}
  \and 
  Vanshika, CSED%
  \thanks{Roll No: \texttt{1024030166}, %
    Email: \texttt{vbansal1\_be24@thapar.edu}}
  \and 
  Nandini Garg, CSED%
  \thanks{Roll No: \texttt{1024030169}, %
    Email: \texttt{ngarg2\_be24@thapar.edu}}
}

\date{\today}

\begin{document}
\maketitle

\section{Higher-order goal%
  \BvrHint{\ldots secondary framing}}
This project aims to help small and mid-sized retail and
wholesale businesses run their day-to-day operations
online through a single system, instead of juggling
separate spreadsheets, ledgers, and point tools. While the
long-term vision is a full-fledged Enterprise Resource
Planning (ERP) suite, the immediate engineering objective
is to translate \emph{unified business operations} into a
measurable, deployable web platform.

\section{Time-to-value %
\BvrHint{\ldots primary heuristic}}
\begin{itemize}[leftmargin=0.7cm]
    \item We will deliver meaningful increments quickly by prototyping the core inventory-to-invoice workflow first and validating it with a mock business within a short iteration window.
    \item We will prioritize fast feedback over perfection by using weekly iterations (and targeting CI-backed automation early) to reduce release friction.
    \item Success is defined by what can be delivered and measured in the first iteration (a working sales-and-inventory loop), then improved based on feedback.
\end{itemize}

\section{Problem Statement}
Retail and wholesale businesses, especially small and
mid-sized ones, typically rely on a patchwork of tools:
paper ledgers, spreadsheets, WhatsApp for orders, and
separate billing software. Common pain points include:
\begin{itemize}[leftmargin=0.7cm]
    \item \textbf{Fragmentation:} inventory, sales, finance, and employee records live in disconnected systems, causing lost context and duplicate data entry.
    \item \textbf{Low visibility:} owners cannot easily see real-time stock levels, cash flow, or sales trends across locations.
    \item \textbf{Manual overhead:} invoicing, payroll, and customer follow-ups are handled manually, consuming staff time and inviting errors.
    \item \textbf{No forecasting:} businesses restock and staff reactively, since there is no systematic demand or cash-flow forecasting.
\end{itemize}

\textbf{Why this matters now (contextual relevance):} As small and mid-sized retailers move online and face thinner margins, even modest automation of inventory, billing, and reporting can materially reduce lost sales, stockouts, and administrative overhead.

\section{Proposed Solution %
\BvrHint{\ldots Software Proposition}}
\subsection{Overview}
Build a \textbf{web-based} ERP platform, \textbf{Vyapar360}, that unifies the core operational and administrative functions a retail or wholesale business needs to run online:
\begin{itemize}[leftmargin=0.7cm]
    \item \textbf{Inventory:} stock tracking, reorder alerts, multi-location and multi-SKU support.
    \item \textbf{Sales:} order capture, point-of-sale style billing, order status tracking.
    \item \textbf{Finance:} ledgers, expense tracking, profit/loss summaries.
    \item \textbf{CRM:} customer records, purchase history, follow-up reminders.
    \item \textbf{Employees \& Payroll:} staff records, attendance, salary computation and disbursal tracking.
    \item \textbf{Invoices:} generation, templating, and GST/tax-ready billing.
    \item \textbf{Analytics:} dashboards for sales, inventory turnover, and financial health.
    \item \textbf{AI chatbot:} conversational assistant for staff/customers to query stock, orders, or FAQs.
    \item \textbf{Forecasting:} lightweight demand and cash-flow forecasting to guide restocking and staffing.
\end{itemize}

\subsection{Core workflow}
\begin{enumerate}[leftmargin=0.7cm]
    \item Staff/owner onboards the business: catalog, initial stock, and employee records.
    \item Sales are recorded through the POS/sales module, automatically decrementing inventory and generating an invoice.
    \item Finance and payroll modules update ledgers and salary computations from sales, expense, and attendance data.
    \item Analytics and forecasting modules surface trends and reorder/staffing suggestions; the chatbot lets staff query the system in natural language.
\end{enumerate}

\subsection{Operational constraints for an educational, course-project setting}
\begin{itemize}[leftmargin=0.7cm]
    \item Keep the solution usable on low-to-mid range devices via responsive web design.
    \item Avoid dependence on advanced research algorithms; focus on workflow, automation, and system correctness, using off-the-shelf forecasting and NLP components rather than novel models.
    \item Design for deployability using common web hosting and managed databases.
\end{itemize}

\section{Solution Approach %
\BvrHint{\ldots Engineering focus}}
This is an engineering course project, so we focus on scalable module design and integration rather than novel algorithm research.

\subsection{Forecasting and chatbot assistance %
\BvrHint{\ldots non-research}}
Forecasting and the chatbot will be built on mature, off-the-shelf approaches:
\begin{itemize}[leftmargin=0.7cm]
    \item Demand forecasting via a simple moving-average or lightweight time-series model (e.g., exponential smoothing) rather than state-of-the-art forecasting research.
    \item AI chatbot via a hosted LLM API constrained to a defined set of intents (stock lookup, order status, FAQs) with retrieval over the business's own data.
    \item Always show forecasts/chatbot answers as suggestions to staff, never as fully automated, unreviewed actions (e.g., no auto-placed purchase orders without confirmation).
\end{itemize}

\subsection{Web architecture %
\BvrHint{\ldots web-based solution}}
A typical 3-tier, modular web architecture:
\begin{itemize}[leftmargin=0.7cm]
    \item \textbf{Frontend:} responsive UI for owners, staff, and customers (dashboards, POS screen, catalog/CRM views).
    \item \textbf{Backend API:} authentication and role-based access, module-specific services (inventory, sales, finance, payroll, CRM), and an integration layer for analytics/forecasting/chatbot.
    \item \textbf{Database:} relational store for transactional modules (inventory, sales, finance, payroll) plus a document/vector store for chatbot retrieval.
\end{itemize}

\subsection{CI/CD and fast delivery enabler}
CI/CD will be used to:
\begin{itemize}[leftmargin=0.7cm]
    \item Automatically build and run tests on every change across modules.
    \item Produce repeatable deployment artifacts to staging.
    \item Enable safe and frequent releases of incremental features (module by module).
\end{itemize}

\section{Evaluation Criterion %
\BvrHint{\ldots measurable and attributable}}
We define success using measurable metrics tied to the core sales-to-invoice workflow.

\subsection{Primary evaluation metric}
\textbf{Time-to-invoice (TTI):} time between a sale being recorded and a corresponding invoice being generated and inventory being updated.
\begin{itemize}[leftmargin=0.7cm]
    \item Target: median TTI $\le$ 5 seconds during pilot window.
    \item Attribution: measured from database timestamps (sale creation vs. invoice/inventory update event).
\end{itemize}

\subsection{Secondary evaluation metrics}
\begin{itemize}[leftmargin=0.7cm]
    \item \textbf{Stock accuracy:} percentage agreement between system-recorded stock and physical stock counts during pilot audits.
    \item \textbf{Forecast usefulness:} percentage of forecast-driven reorder suggestions accepted by staff.
    \item \textbf{Chatbot resolution rate:} percentage of chatbot queries answered correctly without escalation to a human.
    \item \textbf{Reliability:} availability of staging/production for login, billing, and inventory operations during business hours (e.g., $\ge$ 99\% uptime in pilot).
\end{itemize}

\subsection{Pilot validation plan %
\BvrHint{\ldots high-level}}
\begin{itemize}[leftmargin=0.7cm]
    \item Recruit 2--3 small retail/wholesale businesses (or simulate with mock storefronts) and 3--5 staff roles.
    \item Compare metrics before/after within the iteration window by logging baseline estimates via short interviews or existing process observations.
\end{itemize}

\section{Scalability %
\BvrHint{\ldots with foresight}}
The proposition should scale in both theoretical and practical senses.

\begin{enumerate}[leftmargin=0.7cm]
    \item \textbf{Scalable (theoretically):} The system should support growth in number of SKUs, transactions, and concurrent users by:
    \begin{itemize}[leftmargin=0.7cm]
        \item decoupling modules (inventory/sales/finance/CRM/payroll as independently deployable services),
        \item enabling pagination and indexing for common queries (stock lookups, invoice search),
        \item using stateless API design for horizontal scaling.
    \end{itemize}
    
    \item \textbf{Operationally deployable (practically):} The system should run on common web hosting:
    \begin{itemize}[leftmargin=0.7cm]
        \item managed database,
        \item containerized or platform-hosted deployment per module,
        \item CI/CD pipeline for staging/production.
    \end{itemize}
\end{enumerate}

\section{Engine availability heuristic}
A mature set of ``engines'' (tooling/services) is available to implement this as a web-based project:
\begin{itemize}[leftmargin=0.7cm]
    \item Web framework for REST APIs and reactive frontend pages.
    \item Managed relational database for transactional modules; optional document/vector store for chatbot retrieval.
    \item Authentication and role-based access approach (token-based or session-based).
    \item Object storage for product images and invoice PDFs.
    \item Hosted LLM API for the chatbot module.
    \item Charting/BI library for the analytics module.
    \item Test framework for unit/integration tests.
    \item CI runner and staging deployment target.
\end{itemize}
We will use these mature components to focus on module design, integration correctness, and measurement rather than inventing new infrastructure.

\section{Project scope and deliverables %
\BvrHint{\ldots iteration-friendly}}
\subsection{Initial deliverable %
\BvrHint{\ldots first iteration}}
\begin{itemize}[leftmargin=0.7cm]
    \item Inventory module: add/update/view stock with low-stock alerts
    \item Sales module: record a sale, auto-generate an invoice, decrement stock
    \item Basic finance ledger: auto-populated from sales and manual expense entry
    \item Basic logging and timestamps for TTI measurement
    \item CI: build + backend unit tests + linting
    \item CD to staging with a single command or pipeline job
\end{itemize}

\subsection{Subsequent deliverables}
\begin{itemize}[leftmargin=0.7cm]
    \item CRM module: customer records, purchase history, follow-up reminders
    \item Employees \& Payroll module: attendance and salary computation
    \item Analytics dashboard: sales, inventory turnover, and financial summaries
    \item Forecasting module: demand/reorder suggestions
    \item AI chatbot: stock/order/FAQ query assistant
    \item Improved search/filtering and indexed queries for scale readiness
\end{itemize}

\section{Risks and mitigations}
\begin{itemize}[leftmargin=0.7cm]
    \item \textbf{Data privacy / consent:} minimize collected customer and employee data; encrypt sensitive fields (e.g., payroll, payment details); store only what is needed.
    \item \textbf{Staff adoption:} provide an easy UI for billing and stock updates; keep workflow steps minimal; include keyboard-friendly shortcuts for POS use.
    \item \textbf{Forecast/chatbot over-trust:} present forecasts and chatbot answers as suggestions requiring staff confirmation, never as fully automated decisions.
    \item \textbf{Evaluation measurement ambiguity:} instrument timestamps and define clear module boundaries and events before development begins.
    \item \textbf{Operational issues in pilot:} use staging early; ensure CI/CD produces repeatable deployments across all modules.
\end{itemize}

\section{Summary}
This project proposes Vyapar360, a deployable web-based ERP platform that unifies inventory, sales, finance, CRM, employees/payroll, invoicing, analytics, AI chatbot, and forecasting for retail and wholesale businesses. It meets the course criteria by emphasizing time-to-value through rapid, module-by-module iterations, implementing CI/CD to support frequent safe releases, and using measurable evaluation criteria (especially time-to-invoice). It remains focused on engineering workflow, modular architecture, and scalability readiness rather than research-driven novelty, while demonstrating understanding of large-scale, multi-module system design.

\end{document}
